% Comment R3 - 1 

\begin{spacing}{1}
\justifying \color{darkdarkgray} \noindent

Thank you very much for your previous feedback on our paper, which greatly strengthened it. We are glad that you consider we have addressed all your comments adequately.



\end{spacing}


\begin{refcomment}

The paper studies whether reorganizing police work around small, fixed geographic areas improves police effectiveness. The authors exploit the staggered rollout of Plan Cuadrante in Chile and estimate treatment effects using Callaway \& Sant’Anna (2021). They combine household survey data with municipal administrative records on reported crimes and arrests. The paper reports sizeable and policy-relevant effects: a decline in 12-month household victimization, a decline in reported cases, higher arrests, increased trust in police, reduced perceived crime, and lower private security investment. The paper also provides analyses on mechanisms and resources, spillovers, and robustness.

I refereed this paper before for another journal, and I am happy to see that, overall, the paper is strong and much improved. It offers credible evidence that how police are organized matters, not only how many police jurisdictions deploy. I recommend some revisions focused on clarifications and a few additional robustness/interpretation checks.

Below, I reproduce the main points from my prior report and assess whether the current version addressed them.

1.	Mechanism beyond resource expansion. I asked for clearer evidence that specialization—rather than more officers/vehicles—drives the results. The revised manuscript adds two complementary pieces of evidence: (i) early adopters report better police performance in 2005 cross-sections; and (ii) treatment raises perceived police presence and reduces complaints about lack of surveillance. The paper also presents three resource tests: dynamic changes in officers/vehicles after adoption, heterogeneity by actual and predicted resource increases, and no meaningful association between resources and outcomes where Plan Cuadrante had long been in place. The analysis of Seguridad Ciudadana (a later municipal security force program) shows no effects. Together, these additions make a convincing case that specialization—stable assignment to small areas that deepens local knowledge and community contact—explains the gains.

2.	Parallel trends and pre-trends. I asked for a stronger pre-trends discussion. The event-study figures show flat pre-paths, and the text reports joint pre-treatment p-values for all primary outcomes, which supports the identifying assumption.

3.	Administrative measurement and reporting behavior. I asked for a short discussion because higher trust can raise reporting. The manuscript notes that improved trust could bias administrative counts upward, so observed declines likely understate true crime reductions. It also shows consistent declines by crime type and a placebo with non-neighborhood crimes (sexual abuse) that does not change.

4.	External validity. I asked for a short discussion of what features of Chile matter for generalization. The conclusion explicitly flags two contextual features—limited organized crime presence during the study window and relatively high police training—and argues how these shape transferability.

This is careful, policy-relevant work that convincingly shows that geographic specialization of police can reduce crime and improve public confidence. The revisions address virtually all of the concerns in my previous report, and the paper now contains a thorough and transparent robustness and mechanism section.

\label{pt:r3_01}
\end{refcomment}


\begin{response}
    \input{Parts/r3_01.tex}
\end{response}
\stepcounter{AnsweredComments}
